\documentclass{beamer}
\usetheme{Madrid}
\usepackage{amsmath, amssymb, amsthm}
\usepackage{graphicx}
\usepackage{listings}
\usepackage{gensymb}
\usepackage[utf8]{inputenc}
\usepackage{hyperref}
\usepackage{gvv}
\usepackage{tikz}
\lstset{
  language=Python,
  basicstyle=\ttfamily\small,
  keywordstyle=\color{blue},
  stringstyle=\color{orange},
  numbers=left,
  numberstyle=\tiny\color{gray},
  breaklines=true,
  showstringspaces=false
}
\usetikzlibrary{decorations.pathmorphing}

\title{Question-5.2.25}
\author{EE25BTECH11022 - sankeerthan}
\date{}
\begin{document}

\frame{\titlepage}

\begin{frame}
\frametitle{Question}
Solve the following system of linear equations $ax+by =c$, $bx+ay =1+c$
\end{frame}

\begin{frame}
\frametitle{Solution}
Given  
\begin{align}
ax+by &= c \\
bx+ay &=1+c 
\end{align}
The matrix equation for a line is defined as
\begin{align}
    \vec{n}^\top\vec{x}=c
\end{align}
where $\vec{n}$ is the coefficient matrix and $\vec{x}=\myvec{x\\y}$
\begin{align}
       \myvec{a & b}\vec{x}&=c \\ 
       \myvec{b & a}\vec{x}=1+c
\end{align}
\end{frame}
\begin{frame}
\frametitle{Solution}
As a matrix equation
\begin{align}
\myvec{a&b\\b&a}\vec{x}&=\myvec{c \\ 1+c} \\
\myvec{a&b\\b&a}^\top\myvec{a&b\\b&a}\vec{x}&=\myvec{a&b\\b&a}^\top\myvec{c\\1+c} \\
\myvec{a&b\\b&a}\myvec{a&b\\b&a}\vec{x}&=\myvec{a&b\\b&a}\myvec{c\\1+c}   \\
\myvec{a^2+b^2 & 2ab \\2ab & a^2+b^2}\vec{x}&=\myvec{b+\brak{a+b}c \\a+\brak{a+b}c} 
\end{align}
\end{frame}
\begin{frame}
\frametitle{solution} 
\begin{align}
\vec{x}&=\frac{1}{(a^2-b^2)^2}\myvec{a^2+b^2 & -2ab \\-2ab& a^2+b^2}\myvec{b+\brak{a+b}c \\a+\brak{a+b}c} \\
\vec{x}&=\frac{1}{(a^2-b^2)^2} \myvec{\brak{a^2+b^2}\brak{b+\brak{a+b}c}-2ab\brak{a+\brak{a+b}c} \\-2ab\brak{b+\brak{a+b}c}+\brak{a^2+b^2}\brak{a+\brak{a+b}c}} \\
\vec{x}&=\frac{1}{(a^2-b^2)^2}\myvec{b\brak{a^2+b^2-2a^2}+ \brak{\brak{a+b}c}\brak{a^2+b^2-2ab} \\\brak{\brak{a+b}c}\brak{a^2+b^2-2ab}+ a\brak{a^2+b^2-2b^2}} \\
\myvec{x \\y}&=\frac{1}{(a^2-b^2)^2}\myvec{\brak{\brak{a+b}c}\brak{a-b}^2 -b\brak{a^2-b^2} \\\brak{\brak{a+b}c}\brak{a-b}^2 + a\brak{a^2-b^2}} \\
  \end{align}
Hence $x =\frac{c}{a+b}+\frac{b}{b^2-a^2},y=\frac{c}{a+b}+\frac{a}{a^2-b^2}$ 
is the solution for given system of linear equations
\end{frame}
\begin{frame}
\frametitle{graph}
\begin{figure}
    \centering
    \includegraphics[width=0.7\columnwidth]{figs/points.png}
    \caption{}
\end{figure}
\end{frame}
\begin{frame}[fragile]
\frametitle{C-Code}
\begin{lstlisting}[language=C]
// code.c
#include <stdio.h>

// Compute (x, y) for system:
// ax + by = c
// bx + ay = 1 + c

void solve_system(double a, double b, double c, double *x, double *y) {
    double denom = a*a - b*b;
    *x = (a*c - b*(1.0 + c)) / denom;
    *y = (-b*c + a*(1.0 + c)) / denom;
}
\end{lstlisting}
\end{frame}
\begin{frame}[fragile]
\frametitle{C-Code}
\begin{lstlisting}[language=C]
// Fills arrays with points on the line for fixed a, b, c, and range of x values
void generate_points(double a, double b, double c, double *xs, double *ys, int n) {
    double x_min = -5.0, x_max = 5.0;
    for(int i=0; i<n; ++i) {
        xs[i] = x_min + (x_max - x_min)*i/(n-1);
        ys[i] = (c - a*xs[i]) / b;
    }
}
\end{lstlisting}
\end{frame}
\begin{frame}[fragile]
\frametitle{Python-Code}
\begin{lstlisting}
# Code by GVV Sharma
# Modified for Problem Solution
# Released under GNU GPL
# Calculating area enclosed between curves
# call.py
import ctypes
import numpy as np

lib = ctypes.CDLL('./code.so')
\end{lstlisting}
\end{frame}
\begin{frame}[fragile]
\frametitle{Python-Code}
\begin{lstlisting}
# Function signatures
lib.solve_system.argtypes = [ctypes.c_double, ctypes.c_double, ctypes.c_double,
                             ctypes.POINTER(ctypes.c_double), ctypes.POINTER(ctypes.c_double)]
lib.generate_points.argtypes = [ctypes.c_double, ctypes.c_double, ctypes.c_double,
                               np.ctypeslib.ndpointer(dtype=np.float64, flags="C_CONTIGUOUS"),
                               np.ctypeslib.ndpointer(dtype=np.float64, flags="C_CONTIGUOUS"),
                               ctypes.c_int]

def solve_system(a, b, c):
    x = ctypes.c_double()
    y = ctypes.c_double()
    lib.solve_system(a, b, c, ctypes.byref(x), ctypes.byref(y))
    return x.value, y.value
\end{lstlisting}
\end{frame}
\begin{frame}[fragile]
\frametitle{Python-Code}
\begin{lstlisting}
def extract_points(a, b, c, n=20):
    xs = np.zeros(n, dtype=np.float64)
    ys = np.zeros(n, dtype=np.float64)
    lib.generate_points(a, b, c, xs, ys, n)
    return xs, ys

if __name__ == "__main__":
    a, b, c = 2, 1, 5
    xsol, ysol = solve_system(a, b, c)
    print("x =", xsol, "; y =", ysol)
    xs, ys = extract_points(a, b, c, 10)
    print("xs:", xs)
    print("ys:", ys)
\end{lstlisting}
\end{frame}
\begin{frame}[fragile]
\frametitle{Python-Code}
\begin{lstlisting}
# plot.py
import numpy as np
import matplotlib.pyplot as plt
from call import solve_system

# Parameters
a, b, c = 2, 1, 5

# Solve for intersection
x_int, y_int = solve_system(a, b, c)

# Generate x values for plotting
x_vals = np.linspace(-5, 5, 200)

# Calculate y-values for each line
# Line 1: ax + by = c => y = (c - a*x) / b
y_line1 = (c - a * x_vals) / b
\end{lstlisting}
\end{frame}
\begin{frame}[fragile]
\frametitle{Python-Code}
\begin{lstlisting}
# Line 2: bx + ay = 1 + c => y = (1 + c - b*x) / a
y_line2 = (1 + c - b * x_vals) / a

# Plot lines
plt.plot(x_vals, y_line1, label=r'$ax + by = c$')
plt.plot(x_vals, y_line2, label=r'$bx + ay = 1 + c$')

# Plot intersection point
plt.scatter([x_int], [y_int], color='red', marker='X', s=100, label='Intersection Point')

plt.xlabel('x')
plt.ylabel('y')
plt.title('example graph for ax + by = c, bx + ay = 1 + c')
plt.grid(True)
plt.legend()
plt.show()
\end{lstlisting}
\end{frame}

\end{document}